% ==============================================================
% Experiment 2: Zener Diode as a Voltage Regulator
% Meant to be inserted via \input{expt2.tex} in the root document.
% Requires in the root preamble: circuitikz, amsmath, booktabs
% ==============================================================

\chapter{Zener Diode Voltage Regulator}

\section{Aim}
To design, and construct a Zener diode based DC voltage regulator for a given load current and input voltage range, and to study its line regulation
and load regulation characteristics.
\section{Pre-Lab Exercise}
Before performing the experiment, simulate the circuits and verify the designs using LTspice or any other circuit simulation software. 
\section{Apparatus / Components Required}
\begin{center}
	\begin{tabular}{@{}cll@{}}
		\toprule
		S.No. & Component & Specification \\
		\midrule
		1 & Zener diode & $V_Z = 5.6\ \text{V}$, $P_{Z(\max)} = 400\ \text{mW}$ \\
		2 & Resistor $R_S$ & As designed (series/current-limiting resistor) \\
		3 & Resistor (as $R_L$) & $1$--$10\ \text{k}\Omega$ \\
		4 & Regulated DC Power Supply & $0$--$15\ \text{V}$, variable \\
		5 & Digital Multimeter & -- \\
		6 & Breadboard and connecting wires & -- \\
		\bottomrule
	\end{tabular}
\end{center}

\section{Theory}
A Zener diode, when reverse biased beyond its breakdown voltage $V_Z$, maintains
a nearly constant voltage across its terminals over a wide range of reverse
current. This property is exploited to build a simple shunt voltage regulator:
a series resistor $R_S$ limits the current drawn from the unregulated input
$V_{in}$, while the Zener diode, connected in parallel with the load $R_L$,
clamps the output voltage to $V_Z$ as long as it is operated between its
minimum breakdown-sustaining current $I_{Z(\min)}$ and its maximum
rated current $I_{Z(\max)}$ (set by its power rating).

The regulator must continue to hold $V_{out} \approx V_Z$ under two independent
disturbances:
\begin{itemize}
	\item \textbf{Line regulation:} $V_{in}$ varies over its specified range while the load current $I_L$ is held constant.
	\item \textbf{Load regulation:} The load current $I_L$ varies (via $R_L$) while $V_{in}$ is held constant.
\end{itemize}

\section{Design}

\paragraph{Given specifications.}
\begin{align*}
	V_Z &= 5.6\ \text{V} \\
	I_{L(\max)} &= 5\ \text{mA} \\
	V_{in} &: 8\ \text{V to } 15\ \text{V} \quad (V_{in(\min)} = 8\ \text{V},\ V_{in(\max)} = 15\ \text{V}) \\
	P_{Z(\max)} &= 400\ \text{mW}
\end{align*}

\paragraph{Step 1: Maximum Zener current from the power rating.}
\[
I_{Z(\max)} = \frac{P_{Z(\max)}}{V_Z} = \frac{400\ \text{mW}}{5.6\ \text{V}} \approx 71.4\ \text{mA}
\]
Although this is the maximum Zener current, to operate without damage, a much lower current is assumed. Hence, an $I_{Z(max)}$ is assumed to be 50 mA.
\paragraph{Step 2: Choose minimum Zener current.}
To ensure the diode stays in the breakdown region even at worst case
(minimum $V_{in}$, maximum $I_L$), a minimum holding current is chosen as
\[
I_{Z(\min)} = 5\ \text{mA} \quad (\text{a conservative, commonly used design margin})
\]

\paragraph{Step 3: Design $R_S$ for the worst case (minimum $V_{in}$, maximum $I_L$).}
At $V_{in(\min)}$, the source must supply both $I_{L(\max)}$ and $I_{Z(\min)}$:
\[
R_S = \frac{V_{in(\min)} - V_Z}{I_{L(\max)} + I_{Z(\min)}}
= \frac{8 - 5.6}{(5 + 5)\times 10^{-3}}
= \frac{2.4}{0.010} = 240\ \Omega
\]

Choosing the nearest standard (E12) value: $\boxed{R_S = 220\ \Omega}$.

\paragraph{Step 4: Verify at minimum $V_{in}$, full load (checks regulation is maintained).}
\[
I_{T} = \frac{V_{in(\min)} - V_Z}{R_S} = \frac{2.4}{220} \approx 10.9\ \text{mA}
\]
\[
I_Z = I_T - I_{L(\max)} = 10.9 - 5 = 5.9\ \text{mA} \;\; (> I_{Z(\min)} = 5\ \text{mA}) \quad \checkmark
\]

\paragraph{Step 5: Verify at maximum $V_{in}$, no load (checks Zener power rating is not exceeded).}
\[
I_{Z} = \frac{V_{in(\max)} - V_Z}{R_S} = \frac{15 - 5.6}{220} \approx 42.7\ \text{mA}
\]
\[
P_Z = I_Z \times V_Z = 42.7\ \text{mA} \times 5.6\ \text{V} \approx 239\ \text{mW} \;\; (< 400\ \text{mW}) \quad \checkmark
\]

\paragraph{Step 6: Power rating of $R_S$.}
\[
P_{R_S} = \frac{(V_{in(\max)} - V_Z)^2}{R_S} = \frac{(9.4)^2}{220} \approx 0.40\ \text{W}
\]
Use a $220\ \Omega,\ 1\ \text{W}$ resistor (adequate safety margin over the calculated $0.40$~W).
\paragraph{Step 7: Selecting Load Resistance $R_L$.}
The load resistance is a variable resistance, preferably a rheostat. The load current should be variable the range $0--I_{Lmax}$. Hence, the minimum resistance should be at least:
\[
R_{L,min}\frac{V_Z}{I_{Lmax}} = \frac{\SI{5.6}{\volt}}{\SI{5}{\milli\ampere}} = \SI{1.12}{\kilo\ohm}
\]
Hence, choose a $\SI{1}{\kilo\ohm}$ resistance as the fixed load, in series with a larger variable resistance ($\SI{10}{\kilo\ohm})$ to vary the load current. 
\paragraph{Final design summary.}
\begin{center}
	\begin{tabular}{@{}ll@{}}
		\toprule
		Parameter & Value \\
		\midrule
		Zener diode & $5.6\ \text{V}$, $400\ \text{mW}$ \\
		Series resistor $R_S$ & $220\ \Omega$, $1\ \text{W}$ \\
		Load current range & $0$ to $5\ \text{mA}$ \\
		Input voltage range & $8\ \text{V}$ to $15\ \text{V}$ \\
		Regulated output & $\approx 5.6\ \text{V}$ \\
		\bottomrule
	\end{tabular}
\end{center}

\section{Circuit Diagram}
The circuit diagram to be set up is given in Fig.~\ref{fig:zener-regulator}.

\begin{figure}[htbp]
	\centering
	\begin{circuitikz}[american, scale=1.5, line width=0.85pt]
		\draw
		(0,0) to[battery, name=batt, l=$V_{in}$,invert,  *-] (0,3)
		(0,3) to[short] (1,3)
		%(1,3) to[short, -*] (1,1.5)
		(1,3) to[rmeter, name=VMip, t= V,l=$0-20\ \text{V}$, -*] (1,0)
		(1,3) to[R, name=seriesR, l=$R_S$, *-*] (3.5,3)
		(3,3) to[short] (5.3,3)
		(3.5,3) to[zzD, name=zenerD, l=$D_Z$, *-*, invert] (3.5,0)
		(5.3,3) to[rmeter, name=VMop,t= V,l=$0-10\ \text{V}$, -* ] (5.3,0)
		(5.3,3) to[rmeter, t= A,l=$0-10\ \text{mA}$, *-] (7.4,3)
		(7.4,3) to[vR, l=$R_L$,invert, -] (7.4,0)
		(7.4,0) to[short] (5,0)
		(3,0) to[short] (5,0)
		(3,0) to[short] (0,0)
		;
		\draw[thick, ->] ([shift={(-0.45,-0.45)}]batt.center) --([shift={(0.45,0.45)}]batt.center);
		%\draw (5,3) node[right]{$V_{out}$};
		%\draw (0,0) to [rmeter, t= V] (4,3);
		%\draw (0,0) to [rmeter, t= A] (2,3);
		\draw (0,0) node[ground]{};
		\node[below right=0.37cm, align=left] at (zenerD.center)  {$5.6\ \text{V}$\\$400\ \text{mW}$};
		\node[below right=0.37cm, align=left] at (VMip.center) {MC};
		\node[below right=0.37cm, align=left] at (VMop.center) {MC};
		\node[below=0.3cm, align=center] at (seriesR.center)  {$220\ \Omega,\ \text{1 W}$};
	\end{circuitikz}
	\caption{Zener diode voltage regulator circuit}
	\label{fig:zener-regulator}
\end{figure}

\section{Procedure}

\paragraph{(a) Line regulation.}
\begin{enumerate}
	\item Connect the circuit as shown in Figure~\ref{fig:zener-regulator} with $R_L$ set to give $I_L = I_{L(\max)} = 5\ \text{mA}$ (i.e., $R_L = V_Z / I_{L(\max)} \approx 1.12\ \text{k}\Omega$). A rheostat of 2.2 k may be used as the load resistance $R_L$.
	\item Vary $V_{in}$ in steps from $8\ \text{V}$ to $15\ \text{V}$, keeping $R_L$ fixed.
	\item Note $V_{out}$ across the Zener/load for each value of $V_{in}$.
	\item Plot $V_{out}$ versus $V_{in}$.
\end{enumerate}

\paragraph{(b) Load regulation.}
\begin{enumerate}
	\item Fix $V_{in}$ at a convenient value within the design range (e.g., $V_{in} = 12\ \text{V}$).
	\item Vary $R_L$ (using the rheostat) so that $I_L$ varies from $0$ (open circuit) to $I_{L(\max)} = 5\ \text{mA}$.
	\item Note $V_{out}$ for each value of $I_L$.
	\item Plot $V_{out}$ versus $I_L$.
\end{enumerate}

\section{Tabulation}

\begin{table}[htbp]
\begin{minipage}{0.5\linewidth}
	\centering
	\begin{tabular}{@{}cc@{}}
		\toprule
		$V_{in}$ (V) & $V_{out}$ (V) \\
		\midrule
		8  & \\
		9  & \\
		10 & \\
		11 & \\
		12 & \\
		13 & \\
		14 & \\
		15 & \\
		\bottomrule
	\end{tabular}
	\caption{Line regulation ($I_L$ held at $5\ \text{mA}$)}
	\label{tab:lineRegulationObs}
\end{minipage}
\begin{minipage}{0.5\linewidth}
	\centering
	\begin{tabular}{@{}cc@{}}
		\toprule
		$I_L$ (mA) & $V_{out}$ (V) \\
		\midrule
		0   & \\
		1   & \\
		2   & \\
		3   & \\
		4   & \\
		5   & \\
		\bottomrule
	\end{tabular}
	\caption{Load regulation ($V_{in}$ held constant)}
	\label{tab:loadRegulationObs}
\end{minipage}
\end{table}

\section{Model Graph}
Plot two graphs from the observations:
\begin{enumerate}
	\item $V_{out}$ (Y-axis) vs. $V_{in}$ (X-axis) — should be nearly flat once $V_{in}$ exceeds the value needed to bring the Zener into breakdown, confirming line regulation.
	\item $V_{out}$ (Y-axis) vs. $I_L$ (X-axis) — should remain close to $V_Z$ up to $I_{L(\max)}$, with a slight droop as $I_L$ approaches the design limit, confirming load regulation.
\end{enumerate}

\section{Calculations from Observations}
\[
\text{\% Line regulation} = \frac{\Delta V_{out}}{\Delta V_{in}} \times 100\%
\]
\[
\text{\% Load regulation} = \frac{V_{out(\text{no load})} - V_{out(\text{full load})}}{V_{out(\text{full load})}} \times 100\%
\]

\section{Result}
A Zener diode voltage regulator was designed for $V_Z = 5.6\ \text{V}$,
$I_{L(\max)} = 5\ \text{mA}$, over an input range of $8$--$15\ \text{V}$, using
a $400\ \text{mW}$ Zener diode and a series resistor $R_S = 220\ \Omega$. The
circuit was constructed and its line regulation and load regulation
characteristics were verified experimentally. 

The maximum load regulation observed =\hspace{2cm}, at $I_L\ =\ $

The maximum line-regulation observed = \hspace{2cm}, at $V_{in}\ =\ $

\section{Questions}
\begin{enumerate}
	\item Why is the Zener diode operated in reverse breakdown for voltage regulation?
	\item What determines the maximum and minimum permissible Zener current?
	\item Why must $R_S$ be designed using the minimum input voltage and maximum load current condition?
	\item What happens to the regulator if $R_L$ is accidentally open-circuited?
	\item Define line regulation and load regulation. What is an ideal value for each?
	\item How does the power rating of the Zener diode limit the input voltage range for a fixed $R_S$?
\end{enumerate}